% TODO(team): Preserve the negative overall result and censored recovery outcome.
\section{E3: policy comparison}
\begin{frame}{E3: How Repair Policies Balance Playback and Recovery}
% \begin{block}{Repair Policies}
% \begin{itemize}
%     \item \textbf{S:} shared upload bandwidth, no additional repair cap.
%     \item \textbf{F25 / F50:} repair capped at 25\% / 50\% of source capacity.
%     \item \textbf{L:} repair allowance adapts to foreground activity and task age.
% \end{itemize}
% \end{block}

\begin{block}{Experimental Setup}
\begin{itemize}
  \item \textbf{Simultaneous arrivals:} all 6 viewers start within
        the first 0.5 seconds.
  \item \textbf{Staggered arrivals:} 6 viewers start in three groups
        of 2, at approximately 0, 4, and 12 seconds.
  \item \textbf{Controlled comparison:} all policies use the same
        video, initial replica placement, failed Peer, and failure timing.
  \item \textbf{Shared capacity:} each Peer's total upload limit
        stays the same across policies; only the repair allowance changes.
\end{itemize}
\end{block}

\takeaway{Evaluate playback completion, stalls, and verified replica recovery together.}
\note{All policies share the same total source upload limiter. F25 and F50
add repair caps rather than reserving bandwidth. Staggered arrivals are at approximately
0, 0, 4, 4.2, 12, and 12.2 seconds. Each workload-policy combination uses three
paired source-failure contexts. Recovery requires matching inventory receipts
and committed registrations. F50 was selected before inspecting final results.
The task-age floor did not activate during these runs.}
\end{frame}

\begin{frame}{Playback and Recovery Across Repair Policies}
\footnotesize
\renewcommand{\arraystretch}{1.15}
\begin{tabularx}{\textwidth}{
  l *{4}{>{\centering\arraybackslash}X}
}
\toprule
Policy / allowance
& \multicolumn{2}{c}{Avg. stall per viewer (s)}
& \multicolumn{2}{c}{Recovered by 40 s} \\
\cmidrule(lr){2-3}
\cmidrule(lr){4-5}
& Simultaneous & Staggered & Simultaneous & Staggered \\
\midrule
S: no extra cap & 4.46 & 6.18 & 3/3 & 2/3 \\
F25: 25\%      & 4.73 & 6.28 & 3/3 & 1/3 \\
F50: 50\%      & 4.88 & 6.20 & 3/3 & 2/3 \\
L: adaptive    & 4.33 & 6.21 & 3/3 & 1/3 \\
\bottomrule
\end{tabularx}
\par\vspace{2mm}
{\footnotesize
Recovered by 40 s: affected segments regaining majority verified within 40 s after failure.
}
\par\vspace{2mm}
{\footnotesize
L: 25\% when busy; 100\% after 2 s idle. All policies share the same total source cap.
}
\par\vspace{2mm}
Completed sessions in 3 runs: \textbf{7/18} simultaneous; \textbf{9/18} staggered. (all policies)
\takeaway{L showed no consistent advantage over F50 in playback and recovery.}
\note{Stall values include all sessions, including failed sessions, and represent
mean cumulative observed stall per viewer. Recovery timing is measured from the
first fault, with three paired contexts per policy and workload. Mean L minus F50
stall difference is -0.55 seconds for simultaneous arrivals and +0.01 for staggered
arrivals. The former is driven by one context. An L recovery observation at 40.05 s
is close to the threshold relative to the approximately two-second verification
cadence. The task-age floor did not activate in these runs. Recovery means all affected segments have two verified, reachable copies. F25 had one staggered run unrecovered by the observation horizon.}
\end{frame}

\begin{frame}{Lower Playback Stalls Can Delay Replica Recovery}
\small
\textbf{Simultaneous arrivals:} the slowest Peer fails; all six viewers finish.
\par\vspace{3mm}
\begin{tabularx}{\textwidth}{Ycc}
\toprule
Metric & F50 (50\%) & L (adaptive) \\
\midrule
Avg. cumulative stall / viewer & 8.53 s & 6.88 s \\
Verified recovery after failure & 33.99 s & 36.06 s \\
One-copy exposure & 18.82 MiB$\cdot$s & 20.38 MiB$\cdot$s \\
\bottomrule
\end{tabularx}
\par\vspace{3mm}
\takeaway{L reduced stalls by 1.65 s per viewer, but repaired 2.07 s later.}
\note{This is one paired workload/failure context, not an overall advantage.
The task-age floor did not activate. Verified recovery requires two reachable,
hash-valid, registered copies of every affected segment.
Source: compact-study-results.md, E3.}
\end{frame}
